\section{Background and Related Work}
\label{sec:literature}

\subsection{Continual Learning}
\label{subsec:continual_learning}
The introduction described catastrophic forgetting and the stability--plasticity dilemma; here we place the methods that address them. Approaches to forgetting fall into three families \cite{van_de_Ven_2025}. Regularisation-based methods add a penalty that slows changes to parameters identified as important for past tasks, where importance is estimated from a diagonal approximation of the loss curvature \cite{Kirkpatrick_2017}. Architecture-based methods give each task its own parameters, avoiding interference at the cost of a model that grows with the number of tasks. Replay-based methods store a small buffer of past examples and interleave them with new-task data during training. This thesis studies the replay family, and in particular its simplest member, Experience Replay.

\subsection{Experience Replay}
\label{subsec:experience_replay}
Experience Replay (ER) is the simplest replay baseline \cite{chaudhry2019tinyepisodicmemories}. At each step it draws a replay mini-batch from the buffer, combines it with the current-task mini-batch, and takes a single optimiser step on the summed loss of the two. The buffer is populated by reservoir sampling \cite{vitter1985random}, so every sample seen so far occupies a slot with equal probability regardless of arrival order. During $T_0$ the buffer is empty by construction, so ER reduces to plain SGD on the first task, and the stability gap is measurable only from $T_1$ onward.

This simplicity is also why ER remains competitive. \citet{chaudhry2019tinyepisodicmemories} show that jointly training on the current task and a small episodic memory significantly outperforms many purpose-built continual-learning methods. Methods that regularise heavily to protect past tasks tend to give up plasticity, whereas a small replay buffer scales more gracefully to longer task sequences. Consequently, ER serves as a good baseline. Because it does not explicitly manipulate or balance the magnitudes and directions of the gradients, it provides a clean setting to isolate the underlying causes of the stability gap. We therefore adopt it as one of our two core baselines.

A finite buffer is an imperfect estimator of the true past-task distribution. \citet{aljundi2019gradientbasedsampleselection} frame buffer selection as a constraint-reduction problem in which the stored samples approximate the feasibility region defined by the full past data; reservoir sampling is one (non-optimal) policy within that frame. The replay gradient computed on a small buffer therefore differs in direction and magnitude from the gradient on all past data. This estimator gap is present at every step and disappears under no schedule; we report it as a separate, persistent contributor to the stability gap.

\subsection{Path-Finding Methods}
\label{subsec:path_finding}
A second class of replay-based and replay-adjacent methods constrains the SGD update at the task boundary using information about the past tasks. They differ in what information that is. GEM stores past-task gradients and projects the new-task gradient onto the nearest direction whose inner product with every stored gradient is non-negative, so the update does not increase the loss on any past task in the buffer \cite{lopezpaz2022gradientepisodicmemorycontinual}; this uses first-order (gradient) information only. A-GEM replaces the per-task constraints with a single averaged past-task gradient, reducing the projection to one cheap operation \cite{chaudhry2019efficientlifelonglearningagem}. Both act only on the current step, not across the trajectory. EWC instead uses second-order information: it adds a quadratic penalty around the past optimum weighted by the diagonal of the Fisher information, slowing changes to parameters the old task constrained \cite{Kirkpatrick_2017}. NCL also draws on curvature, preconditioning the update by the precision matrix of the past-task posterior to take a natural-gradient step along directions that leave the replay loss roughly unchanged \cite{kao2021naturalcontinuallearningsuccess}. In the framing of this thesis (Section~\ref{subsec:pathfinding_vs_shaping}), all four are \emph{path-finding} methods: they accept the discontinuity at the task boundary and use stored gradient or curvature information to pick a safe trajectory across it. NCL is the second baseline we compare against.

\subsection{The Stability Gap}
\label{subsec:stability_gap}
\citet{delange2023continualevaluationlifelonglearning} introduced continual evaluation and used it to identify the stability gap: the transient accuracy dip on previously seen tasks at the start of new-task training. They observe that the gap is largest in the first few hundred steps of each task, that larger buffers produce smaller gaps, that recovery is partial with the minimum accuracy systematically below the pre-task baseline, and that the gap survives reservoir sampling and moderate learning rates. To explain it, they decompose the update gradient into a plasticity component that lowers the new-task loss and a stability component that preserves the old tasks, and trace the gap to the imbalance between the two at the boundary, where the new-task component dominates. \citet{hess2024complementaryperspectivescontinuallearning} reframe this as a distinction between \emph{what} is optimised and \emph{how}: they show the gap persists even with a perfect approximation of the joint loss, so a better objective alone cannot remove it and the optimisation trajectory must be addressed directly. Subsequent work has shown the gap to be broad rather than incidental: it appears in joint incremental learning of homogeneous tasks \cite{kamath2024expandingscopestabilitygap}, in continual pre-training of large language models \cite{guo2024efficientcontinualpretrainingmitigating}, and in relation to the classification head \cite{lapacz2024exploringstabilitygapcontinual}, and mitigations have been proposed, from temporal ensembles \cite{soutifcormerais2023improvingonlinecontinuallearning} to explicit gap-reduction procedures \cite{harun2024overcomingstabilitygapcontinual}.

\subsection{The Gap Survives the Joint-Loss Oracle}
\label{subsec:kao}
The trajectory argument was made first by \citet{kao2021naturalcontinuallearningsuccess}, who isolate a clean fact about gradient descent through a task transition and illustrate it on a low-dimensional toy problem with two task minima and their combined loss. Remove every replay confound --- no buffer, no sampling noise, no magnitude imbalance --- and give the optimiser exact access to the joint loss $L_{\text{joint}} = L_{T_0} + L_{T_1}$, then run SGD from a stationary point $\theta_0^*$ of $L_{T_0}$. This is the strongest possible replay regime, equivalent to multi-task training with perfect knowledge of the past. Even so, accuracy on $T_0$ dips before recovering. At $\theta_0^*$ the past-task gradient is near zero while the new-task gradient is not, so the first step moves essentially along $-g_1$, where $g_1 := \nabla L_{T_1}(\theta_0^*)$. Writing $H_0 := \nabla^2 L_{T_0}(\theta_0^*)$ for the old-task Hessian, a second-order expansion of the old-task loss gives
\begin{equation}
    \begin{aligned}
        L_{T_0}(\theta_1) - L_{T_0}(\theta_0^*)
         & \approx \tfrac{1}{2}\,\eta^2\, g_1^\top H_0\, g_1 \\
         & > 0,
    \end{aligned}
    \label{eq:kao_taylor}
\end{equation}
whenever $H_0$ has positive curvature in the direction of $g_1$, which it generically does at a strict minimum. The dip is not caused by bad gradients --- the gradient is exact --- but by the \emph{trajectory} the optimiser takes. The conceptual figure in our introduction (Figure~\ref{fig:loss_landscape_visualization}) is drawn in the same spirit; its green curve is precisely this steepest-descent trajectory of the joint loss. This trajectory effect is one of three contributors to the gap, which Section~\ref{sec:account} collects into a single decomposition.

\subsection{Blurry Task Boundaries}
\label{subsec:blurry}
A parallel line of work relaxes the assumption that tasks arrive at sharp boundaries. Standard class- or task-incremental benchmarks present each task as a clean block: all of $T_0$, then all of $T_1$, with an instantaneous switch between them. Several settings replace this with a gradual change. In the \emph{task-free} setting \cite{aljundi2019taskfreecontinuallearning} the learner is never told that a boundary has occurred and the data distribution drifts continuously, so there is no privileged moment at which to consolidate. In the \emph{blurry} setting introduced with Rainbow Memory \cite{bang2021rainbowmemorycontinuallearning} the classes of neighbouring tasks overlap in time: samples of an upcoming class begin to appear before the previous task has ended, so each task shades into the next rather than replacing it. Online formulations push this further, studying blurry, boundary-free streams at the level of individual samples under anytime-inference constraints \cite{koh2022onlinecontinuallearningclass}. Common to all of them is that the transition from one task to the next is spread over many steps instead of one.

This gradual transition is what connects the setting to our intervention, and it makes the blurry-boundary line the closest existing work to the $\lambda$-curriculum. Consider a step on a mini-batch that is an $\alpha$-fraction $T_1$ and a $(1-\alpha)$-fraction $T_0$, with $\alpha(t)$ rising from $0$ to $1$ over the transition window. Its expected gradient is $\alpha\,\nabla L_{T_1}(\theta) + (1-\alpha)\,\nabla L_{T_0}(\theta)$, a weighted sum of the two task gradients. Substituting $\lambda = \alpha/(1-\alpha)$ shows this is, up to an overall scale, the gradient of $L_{T_0} + \lambda L_{T_1}$, the loss the $\lambda$-curriculum descends. A blurry data stream and the $\lambda$-curriculum therefore apply the same fix to the one-sided first step from opposite sides: the former mixes the two task gradients through the \emph{data}, the latter through the \emph{loss weight}. Both keep the past-task gradient present from step one while introducing the new-task gradient gradually.

The two differ in scope. Blurry boundaries achieve the mixing through the data stream and need no buffer, but once $\alpha = 1$ the old-task data, and with it the only retention pressure, is gone; the $\lambda$-curriculum is built on replay and retains an explicit replay signal at $\lambda = 1$. The curriculum is also the cleaner controlled experiment for the stability gap: it changes a single weight in the loss while data sampling, buffer contents, and replay magnitude are held fixed, which isolates the gradual ramp of the new-task gradient as the ingredient responsible for any reduction in the gap. The blurry-boundary results are consistent with ours --- softer transitions forget less --- but, because they vary the data distribution and the loss composition together, they do not separate which of the two is doing the work.

\subsection{Linear Mode Connectivity}
\label{subsec:lmc}

The introduction claimed that a low-loss corridor connects the single-task solution to the joint optimum. Work on \emph{mode connectivity} supports this. \citet{garipov2018losssurfacesmodeconnectivity} and \citet{draxler2019essentiallybarriersneuralnetwork} showed that independently trained minima of a deep network, though separated by a barrier along the straight line between them, are joined by curved paths of near-constant low loss. \citet{frankle2020linearmodeconnectivitylottery} sharpened this to \emph{linear} mode connectivity: networks that share a short period of early training can be joined by a straight low-loss path. For continual learning specifically, \citet{mirzadeh2020linearmodeconnectivitymultitask} found that the multitask solution and the sequentially trained solution often lie in the same basin, linearly connected, when the training regime is stable. This matters for the stability gap in two ways. It means the dip is not forced by the geometry, because a low-loss route from $\theta_A^*$ to $\theta_{\text{joint}}^*$ exists; and it sharpens what an intervention has to do, namely keep the optimiser on that route rather than let it bow away, as the trajectory argument predicts it otherwise will.

Taken together, these results circle the same object from different sides --- the trajectory dip is real and survives an exact objective, softer transitions forget less, and a low-loss corridor to the joint optimum exists --- without naming the mechanism that unifies them. The next section supplies it.